% Minimal retained preamble for lane-p2 cards taken from
% sources/sigma-2010-081-generalized-complex-reduction/pinned-original-source.tex
% (Singular Reduction of Generalized Complex Manifolds, DOI 10.3842/SIGMA.2010.081).
%
% Every declaration below is copied verbatim from the author's own file:
%   lines 9-19  : the author's theorem declarations
%   line  21    : \numberwithin{equation}{section}
%   line  23    : \usepackage[all]{xy}
%   lines 26-36 : the author's own macro definitions
% The published paper uses the publisher class \documentclass{sigma} (chosen in
% the \ifpdf block at lines 1-6), which the assembler does not accept -- only
% article and amsart are supported, see NEW_TASK_AUTHORING section 16.1.  The
% wrapper therefore supplies the article class and the base packages
% amsmath/amssymb/amsthm/graphicx.  Nothing was renamed, redefined, added or
% removed; no macro was invented.  xy is retained because the retained body
% contains the author's own commutative diagram.
\documentclass{article}
\usepackage{amsmath,amssymb,amsthm,graphicx}
\usepackage[all]{xy}

\newtheorem{Theorem}{Theorem}[section]
\newtheorem{Lemma}[Theorem]{Lemma}
%\newtheorem{problem}[Theorem]{Problem}
%\newtheorem{conv}[Theorem]{Convention}
\newtheorem{Proposition}[Theorem]{Proposition}
\newtheorem{Corollary}[Theorem]{Corollary}
{\theoremstyle{definition}
\newtheorem{Remark}[Theorem]{Remark}
\newtheorem{Definition}[Theorem]{Definition}
\newtheorem{Example}[Theorem]{Example}
}

\numberwithin{equation}{section}

\newcommand{\deff}[1]{\textbf{\emph{#1}}}
\newcommand{\func}[3]{#1 \colon #2 \to #3}
\newcommand{\hookfunc}[3]{#1 \colon #2 \hookrightarrow #3}
\newcommand{\eval}[2]{\left. #1 \right|_{#2}}
\newcommand{\pont}[1]{\mathbb{T} {#1}}
\newcommand{\cpont}[1]{\mathbb{T}_{\mathbb{C}} {#1}}
%\newcommand{\tang}{\pmb{\top}}
\newcommand{\tang}{T}
\newcommand{\iprod}[2]{\left\langle {#1} , {#2} \right\rangle}
\newcommand{\iiprod}[2]{\left\langle \! \left\langle {#1} , {#2} \right\rangle \! \right\rangle}
\newcommand{\litlet}{\renewcommand{\labelenumi}{(\alph{enumi})}}
